\section{Remove the hypotheses one at a time}
If completeness is omitted, the map $T(x)=x/2$ on $(0,1)$ is a contraction and a self-map but has no fixed point there: the only real solution is zero, which is excluded.

If strict contraction is replaced by $q=1$, the identity map on $\RR$ has every point fixed, while $T(x)=x+1$ has none. Both preserve distances. The boundary $q=1$ therefore does not merely worsen the numerical rate; it can destroy existence or uniqueness.

If the self-map condition is omitted, $T(x)=x/2+1$ on inputs $[0,1]$ has factor $1/2$ but leaves the interval. Its real fixed point is $2$, outside the intended space. A correct inequality cannot repair a missing domain condition.
